\section{Prime cutoffs and scalar smoothing}\label{supp:sec:cutoffs}
This appendix verifies the free cutoff before specializing it to either
saddle. The distinction matters: the hard route deletes all exceptional
starts, whereas the soft scalar route leaves them in the Stein equation.

\subsection{A uniform weighted prime sum}
\begin{proposition}[Free Rankin calculation and both saddles]
\label{supp:prop:cutoff-calculation}
The bounds and expansions in \cref{prop:cutoffs} hold uniformly for $w$
in a fixed multiplicative neighbourhood of $\sqrt{H\log H}$,
where $H=\log N$.
\end{proposition}
\begin{proof}
Every $Y$-defective integer has the unique form $n=a^2s$ with $s$
squarefree and $Y$-smooth. Therefore, for $0<\zeta<1/2$,
\begin{equation}\label{supp:eq:rankin}
 \frac{A(X,Y)}X\le X^{-\zeta}
                         \prod_{p\le Y}(1+p^{-1+\zeta}).
\end{equation}
Indeed sum $\sqrt{X/s}$ over $s\le X$ and use
$s^{-1/2}\le X^{1/2-\zeta}s^{-1+\zeta}$.
Put $\nu=H/w$, choose the larger solution $u\ge1$ of $e^u/u=\nu$,
and set $\zeta=u/w$. Uniformly in the stated range,
\[
 u\sim\tfrac12\log H,\quad
 \nu\asymp\sqrt{H/\log H},\quad
 \zeta\asymp\sqrt{\log H/H},\quad
 \log H=o(\nu).
\]
For $X/N$ in a fixed compact subset of $(0,\infty)$,
$\zeta\log X=u\nu+o(\nu)$.

We require the prime number theorem inside the weighted Euler product,
not merely an upper bound with an unspecified constant. Let
$R(t)=\pi(t)-\operatorname{Li}(t)$. Given $\epsilon>0$, choose a fixed
lower split point beyond which $|R(t)|\le\epsilon t/\log t$.
Stieltjes integration by parts gives, uniformly in $w$,
\begin{align*}
 \sum_{p\le e^w}p^{-1+\zeta}
 &=\int_2^{e^w}\frac{t^{-1+\zeta}}{\log t}\,dt\nobreak
       +O_\epsilon(1)
       +O\left(\epsilon\frac{e^u}{w}
           +\epsilon\int_2^{e^w}\frac{t^{-1+\zeta}}{\log t}\,dt\right)\\
 &=\operatorname{Ei}(u)+o(\nu).
\end{align*}
For the second equality first let $N\to\infty$ and then
$\epsilon\downarrow0$. The substitution $z=\zeta\log t$ costs only
$O(\log(1/\zeta))=o(\nu)$ at the lower endpoint, and
$\operatorname{Ei}(u)\asymp e^u/u=\nu$. Also
$\sum_p p^{-2+2\zeta}=O(1)$, so logarithms of the Euler factors give
\begin{equation}\label{supp:eq:weighted-primes}
 \sum_{p\le e^w}\log(1+p^{-1+\zeta})
       =\operatorname{Ei}(u)+o(\nu).
\end{equation}
Floors in $Y=\lfloor e^w\rfloor$ alter these estimates by a negligible
quantity. Equations \eqref{supp:eq:rankin}--\eqref{supp:eq:weighted-primes}
yield
\[
 \log(A(X,Y)/X)\le-\nu u+\operatorname{Ei}(u)+o(\nu)
                      =-D(\nu)+o(\nu).
\]
A bad integer contaminates at most $B$ windows; $\log B=o(\nu)$ absorbs
this factor. The graph-edge estimate in the article gives
$E_Y/N^2\le e^{-w+o(\nu)}+N^{-1+o(1)}$.
A window uses at most $CBH/w$ large primes, and each prime meets at most
$CB(N/Y+1)$ starts. Hence
\begin{equation}\label{supp:eq:max-degree}
 \Delta_Y\ll B^2(H/w)(N/e^w+1),\qquad
 \Delta_Y/N\le e^{-w+o(\nu)}+N^{-1+o(1)}.
\end{equation}
This maximum-degree bound applies to the graph on all starts, including
exceptional ones.

For a unified expansion write $V_a=aD(H/V_a)$ for $a=1$ or $2$,
$\nu_a=H/V_a$, and $u_a=u(\nu_a)$. Since $D'=u>0$, the defining
function decreases strictly in $V_a$. At zero its sign is positive;
at $H/e$ its sign is negative for sufficiently large $H$. Thus the root
is unique. The parametric identities are
\[
 \nu_a=e^{u_a}/u_a,\qquad
 V_a=a(e^{u_a}-\operatorname{Ei}(u_a)).
\]
Using $\operatorname{Ei}(u)=e^u u^{-1}(1+u^{-1}+O(u^{-2}))$ gives
\[
 \log H=2u_a-\log u_a+\log a+O(u_a^{-1}),\qquad
 u_a=\tfrac12(\log H+\log\log H-\log(2a))+o(1),
\]
and $V_a^2/H=a(u_a-1+O(u_a^{-1}))$. Substitution gives
\eqref{eq:hard-expansion} and \eqref{eq:soft-expansion}.
\end{proof}
The second-order constants belong to this PNT main-term optimization.
They are not matching lower bounds for the probability distance.

\subsection{A finite soft-exception Stein lemma}
\begin{lemma}[Retention of exceptional indicators]\label{supp:lem:soft-retention}
Let a finite indicator family $(I_\alpha)_{\alpha\in\mathcal I}$ have an
exact dependency graph conditionally on $\mathcal F$. Split
$\mathcal I=\mathcal G\sqcup\mathcal X$ deterministically, and assume
$\EE(I_\alpha\mid\mathcal F)=p$ for every $\alpha\in\mathcal G$.
Put $W=\sum I_\alpha$, $\lambda=|\mathcal I|p$,
$\lambda_{\mathcal X}=|\mathcal X|p$ and
$W_{\mathcal X}=\sum_{\alpha\in\mathcal X}I_\alpha$.
Let $p_\beta(\omega)=\EE(I_\beta\mid\mathcal F)(\omega)$. The mean
conditional Poisson distance is at most a constant times
\begin{align}
 &(1\wedge\lambda^{-1})\left[|\mathcal G|p^2+
     p\sum_{\alpha\in\mathcal G}\sum_{\beta\in B_\alpha\setminus\{\alpha\}}
                                     \EE I_\beta
     +\sum_{\alpha\in\mathcal G}\sum_{\beta\in B_\alpha\setminus\{\alpha\}}
                                     \EE(I_\alpha I_\beta)\right]\notag\\
 &\hspace{10mm}+(1\wedge\lambda^{-1/2})
                          (\lambda_{\mathcal X}+\EE W_{\mathcal X}).
 \label{supp:eq:soft-finite}
\end{align}
The two minimum factors are interpreted as one if $\lambda=0$.
No marginal assumption is made for exceptional neighbours.
\end{lemma}
\begin{proof}
If $\lambda=0$, all good indicators vanish almost surely and
$\dTV(\Law(W\mid\mathcal F),\delta_0)\le
\EE(W_{\mathcal X}\mid\mathcal F)$; averaging proves the claim.
For $\lambda>0$ and a target test set $A\subset\NN$, let $f=f_{\lambda,A}$ solve
$\lambda f(k+1)-kf(k)=\1_A(k)-\Pois(\lambda)(A)$.
The standard factors are
\[
 \|f\|_\infty\le1\wedge\sqrt{2/(e\lambda)},\qquad
 \|\Delta f\|_\infty\le(1-e^{-\lambda})/\lambda
                         \le1\wedge\lambda^{-1}
\]
\cite[(2.15)]{Krokowski2017}; see also
\cite[(2.4)]{Rollin2012SteinFactors}.
For $\alpha\in\mathcal G$ set
$U_\alpha=\sum_{\beta\in B_\alpha\setminus\{\alpha\}}I_\beta$ and
$V_\alpha=W-I_\alpha-U_\alpha$. On each fibre $I_\alpha$ is independent
of $V_\alpha$. Add and subtract
$p\EE[f(V_\alpha+1)\mid\mathcal F]$. Telescoping gives
\begin{align*}
 &|p\EE[f(W+1)\mid\mathcal F]-\EE[I_\alpha f(W)\mid\mathcal F]|\\
 &\quad\le\|\Delta f\|_\infty
 \left[p^2+p\sum_{\beta\in B_\alpha\setminus\{\alpha\}}p_\beta
             +\sum_{\beta\in B_\alpha\setminus\{\alpha\}}
                            \EE(I_\alpha I_\beta\mid\mathcal F)\right].
\end{align*}
The remaining term is bounded directly by
\[
 |\EE[\lambda_{\mathcal X}f(W+1)-W_{\mathcal X}f(W)\mid\mathcal F]|
 \le\|f\|_\infty
                  (\lambda_{\mathcal X}+\EE[W_{\mathcal X}\mid\mathcal F]).
\]
The bound is uniform in the test set, so taking the supremum on each
fibre and then averaging proves \eqref{supp:eq:soft-finite}.
\end{proof}

\subsection{Applying the lemma at a movable cutoff}\label{supp:sec:soft-application}
Take $\mathcal X=D_Y$, $p=2^{-L}$ and $\lambda=Np>0$.
Set $\ell=\log^+\lambda$ and $K=\max(1,\lambda)$.
The elementary inequalities
\[
 (1\wedge\lambda^{-1})\lambda^2\le K,\qquad
 (1\wedge\lambda^{-1})(\lambda^2+\lambda)\le2K,\qquad
 (1\wedge\lambda^{-1/2})\lambda\le\sqrt K
\]
hold on both sides of intensity one.
For neighbour marginals use the maximum degree of the graph on all sites:
\[
 p\sum_{x\in\mathcal G}\sum_{y\in B_x\setminus\{x\}}\EE J_{y,L}
 \le p\Delta_Y\EE Z_{N,L}
 \le p\Delta_Y(\lambda+p\mathcal M_B).
\]
The first-difference factor makes this at most
$K(\Delta_Y/N)(1+N^{-1/2+o(1)})$.
The true marginals of exceptional neighbours have not been replaced by $p$.
For the joint sum, overlapping starts contribute zero and the local bound
treats touching pairs; the separated pairs give
\[
 \sum_{x\in\mathcal G}\sum_{y\in B_x\setminus\{x\}}
       \PP(J_{x,L}J_{y,L}=1)
 \ll p^2\bigl(N^{1+o(1)}+N\Delta_Y+R_2(N,L)\bigr).
\]
Since $R_2/N^2\ll N^{-1/3+\varepsilon/4}(1+\lambda^{-1})$,
its contribution, after the Stein factor, is at most
$K N^{-1/3+\varepsilon/2}$ after enlarging the threshold.
The diagonal and local polynomial costs are smaller.
Finally the exceptional zeroth-order term is at most
\[
 C\sqrt K\left(\frac{\mathcal M_B}{N}+\frac{|D_Y|}{N}\right)
 \le C e^{\ell/2-D(H/w)+o(H/w)}+K N^{-1/3+\varepsilon/2}.
\]
Together with \eqref{supp:eq:max-degree}, these estimates give
\eqref{eq:soft-free-rate} for all intensities.

At $w=\widehat V_N$, let $\nu=H/w$ and choose a common $\eta_N\ge0$,
$\eta_N\to0$, dominating the displayed remainder terms. Put
\[
 q=e^{-(w-\ell)/2+\eta_N\nu},\qquad
 r=e^{\ell-w+\eta_N\nu},\qquad P=N^{-1/3+\varepsilon/2}.
\]
Up to an absolute constant, the error is at most $q+r+KP$.
If $q+N^{-1/3+\varepsilon}\ge1$, the trivial total-variation bound suffices.
Otherwise $q<1$ implies $\ell\le w$ and $r\le q$.
Since $w=o(\log N)$, eventually $K\le e^w\le N^{\varepsilon/2}$, so
$KP\le N^{-1/3+\varepsilon}$. This proves the soft minimum bound,
including $\lambda<1$, at the same cutoff and on the full $\mathcal F_Y$.
All thresholds depend only on the fixed band and the error margin, not on
its length parameter. The numerical constants in a coarser ledger do not
affect the claimed asymptotic bound.

All terms admit their conditional versions: integrate the fibrewise
bound over $C\in\mathcal F_Y$ and divide by $\PP(C)$. A single factor
$e^I$ multiplies the ledger, where $I=-\log\PP(C)$; the leading soft
exponent consequently has sufficient budget $2I+\log^+\lambda<\widehat V$.
This argument applies only to events measurable in the chosen prime
field. It does not turn arbitrary post-selection on a run statistic into
an arithmetic conditioning event.
